\section{Intuitive Picture}

This section gives the informal picture behind the formalism. The aim is not to define the framework rigorously, but to indicate what is being described.

\subsection{State and Change}

A biological system is not treated as a fixed object. It is treated as a process in which states continue to change. Molecules are produced, degraded, transported, and coupled through reactions. Cellular states change over time.

The important point is that change does not stop, while the system can still preserve an overall organization. Biological stability is therefore not static immobility. It is sustained organization under continuing update.

\subsection{Bias and Flow}

State change is not uniform. Some regions produce particular molecules more strongly. Other regions suppress them. These asymmetries give direction to state change. We call this directionality a flow.

The flow arises from production, diffusion, transport, and interaction. It is not a separate substance. It is the net direction and rate of state update.

\subsection{Stable Regions}

When flows exist, not all states behave in the same way. Some states are maintained more easily, while others decay or are displaced. As a result, regions can appear to which the system repeatedly returns.

Such regions are stable in the dynamical sense. They can absorb nearby deviations and preserve a recognizable structure.

\subsection{Appearance of Form}

Development and morphogenesis are interpreted as the appearance of stable configurations through continued state change. A form is not assumed in advance. It is observed when the dynamics settle into a stable arrangement.

Cellular distributions and organ-level patterns are described as such stable arrangements.

\subsection{Without Design Assumptions}

This view does not require purpose, design, or instruction as primitive concepts. It requires only state change, asymmetry, flow, and stable configurations.

\subsection{Summary}

The intuitive picture is that biological systems can be understood through state change, flow, and stable structure. The formal sections describe this picture using information density \(\I\), information flow \(\J\), and induced potential \(\PhiI\).
